% TOP_PROBLEM: {"id": 30006949, "title": "Exact critical exponents of oriented percolation in 1+1 dimensions", "category_id": 19, "category": {"id": 19, "name": "probability", "display_name": "Probability", "description": "Problems involving probability theory, stochastic processes, and random structures.", "slug": "probability", "order_index": 19, "created_at": "2026-07-31T15:26:25.671Z"}, "status": "open"}
% FIRST_POSTED: 2026-09-27
% !TeX program = lualatex
% ATTEMPT_STATUS: unresolved
\documentclass[11pt]{article}
\usepackage[margin=25mm]{geometry}
\usepackage{fontspec}
\setmainfont{Latin Modern Roman}
\usepackage{amsmath,amssymb,mathtools}
\usepackage{unicode-math}
\setmathfont{Latin Modern Math}
\usepackage{xurl,hyperref}
\hypersetup{colorlinks=true,urlcolor=blue}
\setcounter{secnumdepth}{0}
\setlength{\parindent}{0pt}
\setlength{\parskip}{5pt}
\setlength{\emergencystretch}{3em}
\begin{document}
\section*{\texorpdfstring{Exact critical exponents of oriented percolation in 1+1 dimensions}{Exact critical exponents of oriented percolation in 1+1 dimensions}}\label{30006949}
\subsection{Definitions and mathematical statement}
On vertices $(x,t)\in{\mathbb{Z}}\times{\mathbb{Z}}_{\ge0}$ with $x+t$ even, put independent open directed bonds from $(x,t)$ to $(x\pm1,t+1)$, each with probability $p$. Let $p_c$ be the infimum of parameters allowing a positive probability of an infinite directed path from the origin. At $p=p_c$, set
\[
 \tau(x,t)={\mathbb{P}}((0,0)\to(x,t)),\quad
 \theta_t={\mathbb{P}}((0,0)\to{\mathbb{Z}}\times\{t\}),\quad
 \chi_t=\sum_x\tau(x,t),
\]
\[
 \xi_t=\left(\frac{\sum_xx^2\tau(x,t)}{\chi_t}\right)^{1/2}.
\]
Determine the exact critical exponents $\rho,\eta,\nu$ in
\[
 \theta_t=t^{-\rho+o(1)},\qquad
 \chi_t=t^{\eta+o(1)},\qquad \xi_t=t^{\nu+o(1)}.
\]
Existence of the corresponding logarithmic limits is part of an exact exponent determination. The radius is the square root of the ratio of two expectations, not a samplewise radius conditioned on survival. No leading constants are requested.
\par\smallskip\textit{Formulation and concept references:} \href{https://arxiv.org/abs/1709.08291}{[S1]}.

\subsection{Short English statement}
What are the exact survival, population, and spatial-spread exponents of critical nearest-neighbor oriented percolation in one space dimension?
\subsection{Research attempt}
\paragraph{Scope and source audit (27 September 2026).}
The model has independent nearest-neighbor directed bonds, not independent
sites, and only one spatial dimension. The three quantities are exactly
those in the statement. In particular $\xi_t^2$ is a ratio of expected
counts weighted by squared position; replacing it by an expected samplewise
radius conditioned on survival changes the observable.

The full primary paper [S1], arXiv version 4 of 26 February 2018, was
inspected. Its Theorem 1.1 proves the critical hyperscaling relation
$\nu=\eta+2\rho$ provided at least two of the exponents exist. It does
not determine their individual values or prove the existence of all
three limits without that hypothesis. Numerical predictions listed in
its table are not exact constants. The box-crossing input is a substantial
theorem, attributed in that paper to Duminil-Copin--Tassion--Teixeira;
it is not replaced below by a simulation.

The present attack first derives exact finite-time formulas and an
elementary subcritical test. It then reconstructs the easy side of
hyperscaling with explicit integer-time indices, and checks what remains
undetermined even after the known reverse inequality is supplied. No
mean-field exponent from a branching approximation is transferred to this
one-dimensional model.

\paragraph{The cluster process and its elementary inequalities.}
Let $A_t$ be the random set of occupied sites reachable from the origin
at time $t$, and let $Z_t=|A_t|$. There are at most $t+1$ possible sites,
namely $-t,-t+2,\ldots,t$. Therefore
\begin{equation}\label{a457:masssurvival}
 \theta_t\le\chi_t\le(t+1)\theta_t.
\end{equation}
The first inequality is $\mathbf1_{Z_t>0}\le Z_t$, and the second is
$Z_t\le(t+1)\mathbf1_{Z_t>0}$, followed by expectations.

Time slabs use disjoint bonds. Conditional on $A_s$, the expected number
of descendants at time $s+t$ is at most $|A_s|\chi_t$, since counting
descendants separately from each starting site may count some endpoints
more than once. Thus
\begin{equation}\label{a457:submultiplicative}
 \chi_{s+t}\le\chi_s\chi_t.
\end{equation}
On $A_s\ne\varnothing$, select its leftmost point using the past bonds.
Its future survival probability for $t$ more steps is $\theta_t$,
independently of that past. This gives the other-direction inequality
\[
 \theta_{s+t}\ge\theta_s\theta_t.
\]
Neither inequality is an equality in general. The population upper
bound loses information precisely when descendants from different
ancestors merge at the same endpoint.

\paragraph{Lemma 457.1: a finite-time certificate for subcriticality.}
If for some $k\ge1$ and parameter $p$ one has $\chi_k(p)<1$, then
survival decays exponentially and there is no infinite open path.

\emph{Proof.} Write $t=mk+r$, $0\le r<k$. Iterating
\eqref{a457:submultiplicative} gives
$\theta_t\le\chi_t\le\chi_k^m\max_{0\le r<k}\chi_r$.
The right side decreases exponentially in $m$. Hence
$\Pr(A_t\ne\varnothing\text{ for every }t)=\lim_t\theta_t=0$.
\hfill$\square$

It follows also that $\chi_k(p_c)\ge1$ for every $k$. Otherwise
the identity $\chi_k(1)=k+1$ first rules out $p_c=1$, and
finite-time continuity in $p$ would give the same strict inequality at
some $p>p_c$. The lemma would forbid survival there, contrary to the
definition of $p_c$ and monotonicity of survival in $p$. More explicitly,
the set of parameters with positive survival is upward closed and
nonempty, since it contains one; every parameter strictly above its
infimum lies in that set.

Each length-$t$ directed path is open with probability $p^t$, and there
are $2^t$ such paths. Counting paths instead of distinct endpoints gives
$\chi_t\le(2p)^t$. The familiar branching comparison therefore implies
$p_c\ge1/2$. A direct two-step calculation improves this bound without
any asymptotic estimate.

\paragraph{Exact calculations at times one and two.}
At time one,
\[
 \theta_1=2p-p^2,\qquad\chi_1=2p,\qquad\xi_1^2=1\quad(p>0).
\]
At time two the two outer endpoints each have probability $p^2$.
The center can be reached by two bond-disjoint paths, each with
probability $p^2$, so its probability is $2p^2-p^4$. Hence
\begin{equation}\label{a457:time2}
 \chi_2=4p^2-p^4,\qquad
 \sum_xx^2\tau(x,2)=8p^2,\qquad
 \xi_2^2=\frac8{4-p^2}.
\end{equation}
The $-p^4$ term is the first exact correction for merging paths.

For survival, condition on the number of first-step open bonds. If
exactly one is open, survival through the second step has probability
$2p-p^2$; if both are open, it fails only when all four outgoing bonds
are closed. Therefore
\begin{align*}
 \theta_2
 &=2p(1-p)(2p-p^2)+p^2\bigl(1-(1-p)^4\bigr)\\
 &=4p^2-2p^3-4p^4+4p^5-p^6.
\end{align*}
In particular $\chi_2<1$ for
$p<\sqrt{2-\sqrt3}$, and Lemma 457.1 proves
\[
 p_c\ge\sqrt{2-\sqrt3}>\tfrac12.
\]
This modest rigorous lower bound is included as a worked consequence,
not as a claim about the best known threshold estimates.

\paragraph{Exact finite-state recursion, and why it is not an exponent proof.}
Given a reachable set $A$ at time $t$, each site $y$ on the next time
slice has $k_y\in\{0,1,2\}$ incoming bonds from $A$. Its conditional
occupation probability is
\[
 q_y=1-(1-p)^{k_y}.
\]
For distinct sites $y$, these incoming bond sets are disjoint, so their
next-step occupation indicators are independent conditional on $A$.
Thus the transition probability to a particular next set $B$ is exactly
\begin{equation}\label{a457:transfer}
 \Pr(A_{t+1}=B\mid A_t=A)
 =\prod_{y\in B}q_y\prod_{y\notin B}(1-q_y),
\end{equation}
where the product ranges over the finite set of possible children of $A$.
The empty set is absorbing.

Starting from $A_0=\{0\}$, this recurrence computes exact polynomial
or rational values of survival, population, and the second moment at any
fixed finite time. Its number of possible subsets grows exponentially
with the slice width. More importantly, even unlimited finite-time
accuracy does not identify a logarithmic limit without an error estimate
valid uniformly as $t$ grows and $p$ approaches the exact critical point.
A simulation at a nearby decimal parameter can eventually probe the
wrong side of criticality.

\paragraph{Lemma 457.2: a radius bound that respects the observable.}
For every $p>0$ and $t\ge1$,
\begin{equation}\label{a457:radiusbounds}
 1/\sqrt2\le\xi_t\le t.
\end{equation}

\emph{Proof.} The upper bound follows from $|x|\le t$ on the support
of $\tau(\cdot,t)$. For the lower bound, the two possible arrivals at
$x$ at the next time have probabilities $p\tau(x-1,t)$ and
$p\tau(x+1,t)$. Their union has probability at least half their sum,
regardless of the dependence between the parent occupancies. Therefore,
writing $M_t=\sum_xx^2\tau(x,t)$,
\[
 M_{t+1}\ge\frac p2\sum_y\bigl((y-1)^2+(y+1)^2\bigr)\tau(y,t)
             =p(M_t+\chi_t).
\]
The union upper bound gives $\chi_{t+1}\le2p\chi_t$, whence
$\xi_{t+1}^2\ge(\xi_t^2+1)/2\ge1/2$.
Together with $\xi_1=1$, this proves the claim.\hfill$\square$

No samplewise normalization was used. For comparison,
$\mathbb E[Z_t\mid Z_t>0]=\chi_t/\theta_t$ is an exact conditional
population identity, while $\xi_t^2$ weights surviving realizations
according to their population. A configuration with twice as many
occupied points has twice the weight in the underlying expected measure.

\paragraph{The elementary half-time hyperscaling inequality.}
Fix $x,t$, and set $s=\lfloor t/2\rfloor$. A path from $(0,0)$ to
$(x,t)$ requires both survival from the origin to time $s$ and a backward
path from $(x,t)$ to time $s$. These two events use disjoint time slabs.
Time reversal and translation identify the latter probability with
$\theta_{t-s}$. Thus
\begin{equation}\label{a457:pointcap}
 \tau(x,t)\le\theta_s\theta_{t-s}.
\end{equation}
At least three quarters of the mass of $\tau(\cdot,t)$ lies in
$|x|\le2\xi_t$: otherwise its second moment would exceed
$\chi_t\xi_t^2$. There are at most $2\xi_t+1$ allowed parity sites in
that interval. Consequently
\begin{equation}\label{a457:hyperscale}
 \chi_t\le\tfrac43(2\xi_t+1)\theta_{\lfloor t/2\rfloor}
                                      \theta_{\lceil t/2\rceil}.
\end{equation}
This is the elementary mechanism behind the upper inequality in [S1],
with the parity count retained. It is a reconstruction of a known bound,
not a new hyperscaling theorem.

If all three logarithmic exponents exist, the radius lower bound ensures
$\log(2\xi_t+1)/\log t\to\nu$. Taking logarithms in
\eqref{a457:hyperscale} then yields
\[
 \eta+2\rho\le\nu.
\]
Also $\chi_t(p_c)\ge1$, $\theta_t\le1$, and
\eqref{a457:masssurvival}--\eqref{a457:radiusbounds} give the elementary
constraints $\eta\ge0$, $\rho\ge0$, $0\le\nu\le1$, and
$\eta+\rho\le1$. These are constraints conditional on the relevant
limits, not numerical determinations of them.

\paragraph{What the deep reverse inequality supplies, and what it does not.}
Theorem 1.1(ii) of [S1] supplies a constant $c>0$ such that at criticality
\[
 \chi_t\ge c\,\xi_t\theta_t^2.
\]
Its proof uses critical box crossings, rather than the elementary
branching comparison. We use this only as an attributed theorem. Combined
with the upper estimate, it gives the stated hyperscaling equality under
the existence hypotheses in [S1]. It removes one degree of freedom among
the exponents; it does not choose their values. For example, the positive
triples $(\rho,\eta,\nu)=(1/8,1/4,1/2)$ and
$(1/6,1/3,2/3)$ both satisfy the equality and all the elementary bounds
displayed above. Neither is asserted to be the actual triple.

Even a pointwise relation among three sequences need not create their
logarithmic limits. For a diagnostic model at real $t\ge e$, take
\[
 a(t)=\tfrac15+\tfrac1{50}\sin(\log\log t),\qquad
 \Theta(t)=t^{-a(t)},\quad X(t)=1,\quad R(t)=t^{2a(t)}.
\]
Then $X=R\Theta^2$ exactly, $\Theta$ decreases because
$a(t)+\tfrac1{50}\cos(\log\log t)>0$, and the analogous coarse
population and radius bounds hold. But the logarithmic powers of
$\Theta$ and $R$ oscillate and have no limits. These are not claimed to
be percolation observables or to satisfy every cluster inequality. They
refute the shortcut from the displayed scaling relation and coarse
bounds alone to existence of the exponents.

\paragraph{Why branching and finite slopes do not settle the target.}
The path-count process ignores mergers. Already \eqref{a457:time2}
shows that its mean differs from the cluster population. At large times
in one spatial dimension, many directed paths share endpoints and edges;
independence of offspring cannot be assumed. A mean-field exponent for
a branching model therefore requires an additional comparison sharp
enough on the polynomial scale, which is not proved here.

Likewise a measured slope such as
$-\log(\theta_{2t}/\theta_t)/\log2$ is a finite-scale statistic.
It can be stable over a long interval without converging as $t\to\infty$,
and its value at an approximate $p_c$ has an additional parameter error.
Finite transfer-matrix calculations provide exact data at their chosen
time and parameter, but neither interchange of limits is justified by
that exactness. The first unresolved analytic step is a quantitative
control of scale changes at criticality strong enough to establish
logarithmic limits and then identify them.

\paragraph{Verification and outcome.}
The companion script implements \eqref{a457:transfer} with rational
probabilities and checks the one- and two-step formulas, the survival and
population inequalities, half-time pointwise caps, and moment definitions
at several finite times and rational parameters. No decimal estimate of
$p_c$ is treated as its exact value. The computations are therefore checks
of the finite lemmas rather than fitted exponent evidence.

\textbf{Outcome: UNRESOLVED.} The attempt supplies exact finite-time
formulas, an explicit subcritical certificate, moment inequalities, and
an audited reconstruction of the easy hyperscaling bound. The known
reverse bound is attributed. Existence and exact values of the selected
three critical exponents remain unproved by this work.

\subsection{Sources}
\begin{itemize}
\item[S1]A. Sakai, Hyperscaling for oriented percolation in 1+1 space-time dimensions, Journal of Statistical Physics 171 (2018), 462-469.
\url{https://arxiv.org/abs/1709.08291}.
\textit{Formulation source. Consulted 24 September 2026: The full text defines the survival, population, and gyration-radius quantities.}
\item[E1]Akira Sakai, full primary manuscript corresponding to [S1],
arXiv:1709.08291v4, 26 February 2018, Theorem 1.1 and Section 2.
\url{https://arxiv.org/pdf/1709.08291}.
\textit{Directly inspected 27 September 2026. The exact hyperscaling
relation is conditional on exponent existence as stated there; the
individual predicted values are not used as established constants.}

\end{itemize}
\end{document}
